%% =====================================================================
%%  T-14-05  The Bait
%%  -------------------------------------------------------------------
%%  One idea per file. Core is mandatory. Slots in canonical order.
%%  Max 700 words. No layout commands. No filler. New source material
%%  for this entry goes in sources/B3/T-14-05.tex -- never by
%%  rewriting this file (docs/RULES.md R0.1).
%%  =====================================================================
%% @kind: topic
%% @id: T-14-05
%% @title: The Bait
%% @subtitle: Make them come to you --- the one who moves first holds the worse cards
%% @chapter: c14
%% @order: 50
%% @status: stable
%% @sources: B3
%% @updated: 2026-09-19

\topic{T-14-05}{The Bait}{Make them come to you --- the one who moves first holds the worse cards}

\begin{core}
When you force a person to act, you are the one exposed; when you make them come to you, they abandon their own plans in the process and you hold the arms. The move is the bait: place a lure on your ground --- a gain, a curiosity, an unfinished thing --- and let the counterpart's own desire do the travelling. Negotiation from the stationary position is negotiation on chosen terrain.
\end{core}
\begin{mechanism}
B3's judgment: the better position always belongs to the party that does not move first. The mechanism is the asymmetry of the mover: the party that travels spends time, anticipation and initiative, and arrives already invested --- and a person who has invested in reaching the table unconsciously grants the table's owner standing. The bait works on the open loop: a partial, tantalising offer (the fabulous gain dangled, then half-withdrawn) installs an unfinished task the mind keeps working; the counterpart returns to close it, on your schedule, having privately re-priced the prize upward during the wait. This is the aggressive half of this chapter's family: T-14-06 controls the options once they arrive; the bait decides whether and when they arrive at all. The skills version is the withdrawal chapter's law (T-10-05) read as a weapon --- scarcity of access is the standing bait.
\end{mechanism}
\begin{conditions}
Strongest where the counterpart wants something you are not obliged to release --- a sale, a concession, a job, information --- and where they can reach you. It fails against the genuinely indifferent (no desire, no travelling), misfires on time-symmetric standoffs (both sides equally able to wait), and curdles where the bait is false advertising: the lure must contain some real gain, or the return trip ends in a burned door.
\end{conditions}
\begin{application}
  \item Price their want before setting anything. A bait nobody wants is decoration.
  \item Place the offer on your ground: your meeting, your process, your calendar.
  \item Give the lure an open loop --- partially disclosed terms, a detail withheld, a trial that stops just short.
  \item Then stand still. Every chase is a payment against your own position; audit your own chasing weekly.
  \item When they arrive, be slow to close. The travelling party must be allowed to buy; hurrying them re-prices you as the needy one (T-11-02).
\end{application}
\begin{feedback}
  \item Working: the counterpart initiates, returns, and asks what it would take. Terrain is yours.
  \item Working: your terms stop being negotiated and start being requested.
  \item Failing: you are checking whether they read your message. You are the fish.
  \item Failing: they name the bait --- \emph{you're dangling that to draw me in}. The loop closed backwards, with you in it.
\end{feedback}
\begin{failure}
Baited ground poisons repeat trade: a counterpart who realises the meeting was always a trap brings that learning to every future table, and deservedly. The tactic also selects for greed and punishes straight dealing in the operator's own habits --- the person who always fishes starts to regard every counterpart as a catch, including the ones who were partners. And against a player who simply refuses to travel, the bait sits there ageing in public view.
\end{failure}
\begin{countermeasures}
  \item Before travelling to any table, ask who is baiting whom: who set the ground, who named the time, who benefits from your wanting this?
  \item Feel your own open loops --- the unanswered offer, the unfinished deal you keep drifting back to. Name the lure and the pull gets a leash (T-10-02).
  \item Test scarcity claims by walking. True scarcity survives a walk-away; staged scarcity follows you.
  \item In your own selling, use bait you would be content to be caught with: the ethical version of this entry is making your offer genuinely worth travelling to.
\end{countermeasures}

\sources{B3}
\seesources
\seealso{T-14-06, T-10-05, T-11-02}
